\section{量子光场}
\label{chap:quantum-light-states}

把 Maxwell 正常模量子化以后，每个模式都成为一个量子谐振子。接下来要区分的不是不同的场方程，而是同一组模式可以处于哪些量子态，以及光子数涨落和相位噪声怎样在实验中显现。

\subsection{Fock、相干态与一般线性驱动}
对任一固定模式，记 \(a\equiv a_{\mathbf k\lambda}\)。其对易关系

\[
[a,a^\dagger]=1
\]

定义数算符 \(N=a^\dagger a\)。真空满足 \(a|0\rangle=0\)，归一化数态为

\[
|n\rangle
=\frac{(a^\dagger)^n}{\sqrt{n!}}|0\rangle,
\]

并满足

\[
a|n\rangle=\sqrt n|n-1\rangle,
\qquad
a^\dagger|n\rangle=\sqrt{n+1}|n+1\rangle.
\]

Glauber 光探测理论中的一阶关联函数

\[
G^{(1)}(1,2)
=\langle\hat E^{(-)}(1)\hat E^{(+)}(2)\rangle
\]

之所以采用正规 ordering，正是因为理想吸收型探测器先由 \(E^{(+)}\) 湮灭一个光量子，再由共轭振幅组成概率。

Fock 态有确定光子数但无确定相位，是最"量子"的光场。实验上更常见的是激光一类具有确定相位和近经典振幅的光场，其量子态由相干态描述。相干态由湮灭算符本征方程定义：

\[
a|\alpha\rangle=\alpha|\alpha\rangle.
\]

令 \(|\alpha\rangle=\sum_nc_n|n\rangle\)，递推关系

\[
c_{n+1}\sqrt{n+1}=\alpha c_n
\]

给出

\[
c_n=\frac{\alpha^n}{\sqrt{n!}}c_0.
\]

归一化后

\[
|\alpha\rangle
=\ee^{-|\alpha|^2/2}
\sum_{n=0}^{\infty}
\frac{\alpha^n}{\sqrt{n!}}|n\rangle.
\]

光子数服从 Poisson 分布

\[
P_n=\ee^{-|\alpha|^2}\frac{|\alpha|^{2n}}{n!},
\qquad
\langle N\rangle=|\alpha|^2.
\]

\begin{proof}[解]
递推关系。由 $a|\alpha\rangle=\alpha|\alpha\rangle$ 和 $|\alpha\rangle=\sum_n c_n|n\rangle$，左乘 $\langle n|$：
\begin{equation*}
\sum_m c_m\langle n|a|m\rangle=\alpha c_n.
\end{equation*}
由 $a|m\rangle=\sqrt{m}|m-1\rangle$，只有 $m=n+1$ 项贡献：$c_{n+1}\sqrt{n+1}=\alpha c_n$。递推解 $c_n=\alpha^n c_0/\sqrt{n!}$。

归一化。由 $\langle\alpha|\alpha\rangle=1$：
\begin{equation*}
1=\sum_n|c_n|^2=|c_0|^2\sum_n\frac{|\alpha|^{2n}}{n!}=|c_0|^2\ee^{|\alpha|^2},
\end{equation*}
故 $c_0=\ee^{-|\alpha|^2/2}$（取相位为零），得 $|\alpha\rangle=\ee^{-|\alpha|^2/2}\sum_n\frac{\alpha^n}{\sqrt{n!}}|n\rangle$。

光子数分布 $P_n=|\langle n|\alpha\rangle|^2=\ee^{-|\alpha|^2}|\alpha|^{2n}/n!$ 是参数为 $\lambda=|\alpha|^2$ 的 Poisson 分布，其平均光子数为
\begin{equation*}
\langle N\rangle=\langle\alpha|a^\dagger a|\alpha\rangle=|\alpha|^2\langle\alpha|\alpha\rangle=|\alpha|^2,
\end{equation*}
方差 $\langle N^2\rangle-\langle N\rangle^2=|\alpha|^2$（Poisson 分布的特征：方差等于均值）。

正交分量统计。由 $a|\alpha\rangle=\alpha|\alpha\rangle$ 和 $Q=(a+a^\dagger)/\sqrt{2}$，
\begin{equation*}
\langle Q\rangle=\frac{\langle\alpha|a+a^\dagger|\alpha\rangle}{\sqrt2}
=\frac{\alpha+\alpha^*}{\sqrt2}
=\sqrt2\,\operatorname{Re}\alpha,
\end{equation*}
同理 $\langle P\rangle=\sqrt2\,\operatorname{Im}\alpha$。方差由 $(\Delta Q)^2=\langle Q^2\rangle-\langle Q\rangle^2$：
\begin{equation*}
Q^2=\frac{a^2+a^{\dagger2}+aa^\dagger+a^\dagger a}{2}
=\frac{a^2+a^{\dagger2}+2a^\dagger a+1}{2},
\end{equation*}
故 $\langle Q^2\rangle=(\alpha^2+\alpha^{*2}+2|\alpha|^2+1)/2$，而 $\langle Q\rangle^2=(\alpha+\alpha^*)^2/2=(\alpha^2+\alpha^{*2}+2|\alpha|^2)/2$，相减得 $(\Delta Q)^2=1/2$。同理 $(\Delta P)^2=1/2$。不确定乘积 $\Delta Q\cdot\Delta P=1/2$ 恰为 Heisenberg 下界——相干态是最小不确定态。
\end{proof}

定义无量纲正交分量

\[
Q=\frac{a+a^\dagger}{\sqrt2},
\qquad
P=\frac{a-a^\dagger}{\ii\sqrt2},
\]

则

\[
\langle Q\rangle=\sqrt2\operatorname{Re}\alpha,
\qquad
\langle P\rangle=\sqrt2\operatorname{Im}\alpha,
\]

并有 \((\Delta Q)^2=(\Delta P)^2=1/2\)。相干态因此在相空间中保持真空噪声宽度，只让波包中心按经典轨道运动。

相干态不仅是谐振子基态的平移，还可以由经典源驱动谐振子从真空动态产生。以下用 Magnus 展开证明：任意线性驱动（不必恒定或共振）的初态为真空时，末态必为相干态。

考虑单模受经典源驱动：

\[
\hat H
=\hbar\omega(a^\dagger a+1/2)
+\ii\hbar[f(t)a^\dagger-f^*(t)a].
\]

进入自由相互作用表象以后，记相应驱动系数为 \(f_I(t)\)。不同时间的相互作用满足

\[
[\hat H_I(t),\hat H_I(t')]
=\text{c-number}.
\]

因此 Magnus 展开在二阶以后终止。令

\[
G(t)
\equiv-\frac{\ii}{\hbar}H_I(t)
=f_I(t)a^\dagger-f_I^*(t)a.
\]

第一阶 Magnus 项为

\[
\Omega_1
=\int_{t_0}^{t}G(t_1)\dd t_1
=\alpha a^\dagger-\alpha^*a,
\qquad
\alpha(t)=\int_{t_0}^{t}f_I(t_1)\dd t_1.
\]

交换子

\[
[G(t_1),G(t_2)]
=f_I(t_1)f_I^*(t_2)-f_I^*(t_1)f_I(t_2)
\]

是纯虚 c 数，因此所有三重及更高嵌套交换子严格为零。第二阶 Magnus 项只是整体相位：

\[
\Omega_2
=\frac12\int_{t_0}^{t}\dd t_1
\int_{t_0}^{t_1}\dd t_2\,[G(t_1),G(t_2)]
\equiv\ii\varphi(t),
\]

其中

\[
\varphi(t)
=\operatorname{Im}
\int_{t_0}^{t}\dd t_1
\int_{t_0}^{t_1}\dd t_2\,
 f_I(t_1)f_I^*(t_2).
\]

于是时间序指数精确写成

\begin{equation}
U_I(t,t_0)
=\ee^{\ii\varphi(t)}D[\alpha(t)].
\label{eq:driven-oscillator-exact}
\end{equation}

这里的几何-like 相来自不同时刻位移在相空间中的非交换性；虽然它不影响单个相干态的光子数，但在多路径干涉中可以产生可观测相位。由真空出发立刻得到

\[
|\psi_I(t)\rangle
=\ee^{\ii\varphi(t)}|\alpha(t)\rangle.
\]

因此``经典相干源产生相干态''并不依赖恒定驱动或严格共振；那只是上式的特殊情形。

线性驱动保持相干态空间不变，但许多物理过程（参量下转换、频率突变、Bogoliubov 凝聚）的哈密顿量含 $a^{\dagger2}$ 或 $a^2$ 项，属于二次型。这类哈密顿量产生压缩态——另一类从真空出发可动态生成的高斯态。

\subsection{二次哈密顿量、Bogoliubov 变换与高斯态}
一般单模二次哈密顿量可写为

\[
\hat H
=\hbar\omega a^\dagger a
+\frac{\hbar}{2}
\left(\kappa_{\rm sq} a^{\dagger2}+\kappa_{\rm sq}^*a^2\right)
+\text{linear terms}.
\]

其中 \(\kappa_{\rm sq}\in\mathbb C\) 是具有角频率量纲的二光子（parametric）耦合系数；下标 ``sq'' 用来与强场隧穿尺度和 Ray 不对称参数区分。

Heisenberg 方程对 \((a,a^\dagger)^T\) 是线性的，因此任意演化都具有 Bogoliubov 形式

\[
a' = ua+va^\dagger,
\qquad
|u|^2-|v|^2=1.
\]

这个约束正是保持 \([a',a'^\dagger]=1\) 的条件。

\begin{proof}[解]
含参量下转换哈密顿量 \(\hat H=\hbar\omega a^\dagger a+\frac{\hbar}{2}(\kappa a^{\dagger2}+\kappa^*a^2)+\cdots\) 中，\(a^{\dagger2}\) 项成对产生光子。Heisenberg 方程
\begin{equation*}
\frac{\dd a}{\dd t}=\frac{1}{\ii\hbar}[a,\hat H]
=-\ii\omega a-\ii\kappa^* a^\dagger+\cdots
\end{equation*}
解为线性变换 \(a(t)=u(t)a(0)+v(t)a^\dagger(0)+\text{c-number}\)。

正则对易关系保持要求 \([a(t),a^\dagger(t)]=1\)：
\begin{equation*}
[a(t),a^\dagger(t)]=[ua+va^\dagger,u^*a^\dagger+v^*a].
\end{equation*}
展开利用 \([a,a^\dagger]=1\)、\([a,a]=[a^\dagger,a^\dagger]=0\)：
\begin{equation*}
[a(t),a^\dagger(t)]=u v^*[a,a^\dagger]+v u^*[a^\dagger,a]=(uv^*-vu^*)[a,a^\dagger]=|u|^2-|v|^2,
\end{equation*}
最后一步利用 \(uv^*-vu^*=|u|^2-|v|^2\)（对复数 \(u,v\)）。故约束为
\begin{equation*}
|u|^2-|v|^2=1.
\end{equation*}
在压缩相位 \(\phi=0\) 时，Bogoliubov 变换为 \(S^\dagger a S=a\cosh r-a^\dagger\sinh r\)（由\eqnrefs{eq:squeezing-bogoliubov-transform}，\(u=\cosh r\)，\(v=-\sinh r\)，满足 \(|u|^2-|v|^2=\cosh^2 r-\sinh^2 r=1\)）。由 \(Q=(a+a^\dagger)/\sqrt{2}\)：
\begin{align*}
S^\dagger Q S
&=\frac{S^\dagger a S+S^\dagger a^\dagger S}{\sqrt2}
\\
&=\frac{a\cosh r-a^\dagger\sinh r}{\sqrt2}
  +\frac{a^\dagger\cosh r-a\sinh r}{\sqrt2}
\\
&=\frac{(a+a^\dagger)(\cosh r-\sinh r)}{\sqrt2}
\\
&=Q\,e^{-r},
\end{align*}
其中用了 \(\cosh r-\sinh r=e^{-r}\)。同理 \(S^\dagger P S=P\,e^{r}\)。对压缩真空 \(|0;\zeta\rangle=S(\zeta)|0\rangle\)，\(Q\) 方差为
\begin{equation*}
(\Delta Q)^2=\langle 0;\zeta|Q^2|0;\zeta\rangle=\langle 0|S^\dagger Q^2 S|0\rangle=\langle 0|e^{-2r}Q^2|0\rangle=\frac{1}{2}e^{-2r},
\end{equation*}
因为真空 \(\langle 0|Q^2|0\rangle=1/2\)。同理 \((\Delta P)^2=\frac{1}{2}e^{2r}\)。不确定乘积 \(\Delta Q\cdot\Delta P=1/2\) 保持 Heisenberg 下界。

\end{proof}

把

\[
u=\ee^{-\ii\theta}\cosh r,
\qquad
v=-\ee^{-\ii(\theta-\phi)}\sinh r
\]

参数化，就得到相位旋转与压缩。

压缩算符取

\[
S(\zeta)
=\exp\!\left[
\frac12(\zeta^*a^2-\zeta a^{\dagger2})
\right],
\qquad \zeta=re^{\ii\phi}.
\]

在 \(\phi=0\) 时

\[
S^\dagger QS=\ee^{-r}Q,
\qquad
S^\dagger PS=\ee^{+r}P,
\]

所以

\[
(\Delta Q)^2=\frac12\ee^{-2r},
\qquad
(\Delta P)^2=\frac12\ee^{2r}.
\]

\begin{proof}[解]
定义一参数算符
\[
S_s=\exp\left[\frac{s}{2}(\zeta^*a^2-\zeta a^{\dagger2})\right],
\qquad 0\le s\le1,
\]
并记
\[
A(s)=S_s^\dagger aS_s,
\qquad
B(s)=S_s^\dagger a^\dagger S_s.
\]
令
\[
G=\frac12(\zeta^*a^2-\zeta a^{\dagger2}),
\]
则
\[
[a,G]=-\zeta a^\dagger,
\qquad
[a^\dagger,G]=-\zeta^*a.
\]
因此
\[
\frac{\dd A}{\dd s}=-\zeta B,
\qquad
\frac{\dd B}{\dd s}=-\zeta^*A.
\]
写 \(\zeta=r\ee^{\ii\phi}\)，可得
\[
\frac{\dd^2A}{\dd s^2}=r^2A,
\qquad
A(0)=a,
\qquad
A'(0)=-r\ee^{\ii\phi}a^\dagger.
\]
故在 \(s=1\) 时
\begin{equation}
S^\dagger(\zeta)aS(\zeta)
=a\cosh r-\ee^{\ii\phi}a^\dagger\sinh r.
\label{eq:squeezing-bogoliubov-transform}
\end{equation}
这一步同时证明了压缩是保持正则对易子的双曲旋转，而不是对 \(a\) 的普通酉相位旋转。
\end{proof}

对压缩 vacuum \(|0;\zeta\rangle=S(\zeta)|0\rangle\)，由\eqnrefs{eq:squeezing-bogoliubov-transform}可知平均光子数

\[
\langle N\rangle=\sinh^2r.
\]

数态展开只含偶数光子：

\[
|0;\zeta\rangle
=\frac1{\sqrt{\cosh r}}
\sum_{n=0}^{\infty}
\frac{\sqrt{(2n)!}}{2^n n!}
\left[-\ee^{\ii\phi}\tanh r\right]^n
|2n\rangle.
\]

所以压缩并不是“把真空噪声椭圆压扁”的纯几何图像；从原 Fock 基看，它等价于成对地产生偶数个量子。二次哈密顿量中的 \(a^{\dagger2}\) 和 \(a^2\) 正是这种 pair 产生/湮灭结构。

压缩态的生成机制不止参量下转换一种。谐振子频率的突然变化也会产生压缩——这是一个纯量子效应，经典力学中不存在类比。

若谐振子频率从 \(\omega_0\) 突然变到 \(\omega_1\)，突变时间 \(\tau_q\) 满足 \(\tau_q\ll\omega_0^{-1},\omega_1^{-1}\)，则在突变瞬间 \(Q,P\) 来不及演化而保持连续。用两种频率各自定义的升降算符消去 \(Q,P\)，得到

\[
a_1=ua_0+va_0^\dagger,
\]

\[
u=\frac12\left(
\sqrt{\frac{\omega_1}{\omega_0}}
+\sqrt{\frac{\omega_0}{\omega_1}}
\right),
\qquad
v=\frac12\left(
\sqrt{\frac{\omega_1}{\omega_0}}
-\sqrt{\frac{\omega_0}{\omega_1}}
\right).
\]

它自动满足 \(|u|^2-|v|^2=1\)。若采用 \(u=\cosh r,v=\sinh r\) 的实参数约定，则

\[
r=\frac12\ln\frac{\omega_1}{\omega_0}.
\]

``突然近似''的控制条件已经由 \(\tau_q\omega_{0,1}\ll1\) 明确给出；相反，若频率变化极慢，绝热定理使系统跟随瞬时本征态而不会产生同样的压缩激发。

\subsection{量子光场的实验表征与典型参数}\label{sec:quantum-optics-experiments}

不同量子态的区别最终要由光子计数和场正交分量的统计量来辨认。二阶关联、Fano 因子和压缩噪声分别测量聚束、数涨落与低于真空的正交分量噪声，并对应不同的实验制备方法。

Glauber \(n\) 阶关联函数中，二阶零时延关联
\[
g^{(2)}(0)=\frac{\langle a^\dagger a^\dagger aa\rangle}{\langle a^\dagger a\rangle^2}
=\frac{\langle N(N-1)\rangle}{\langle N\rangle^2}
\]
是区分光场类型的核心可观测量。HBT (Hanbury Brown--Twiss) 实验用 50/50 分束器后两路探测器符合计数直接测量 \(g^{(2)}(0)\)。

\begin{center}
\begin{tabular}{lccc}
\hline
光场类型 & \(g^{(2)}(0)\) & \(\Delta N^2/\langle N\rangle\) & 实验方案 \\
\hline
Fock 态 \(|n\rangle\) & \(1-\frac{1}{n}\) & \(0\) & 单分子、量子点 \\
相干态 \(|\alpha\rangle\) & \(1\) & \(1\) (Poisson) & 理想激光 \\
热光场 & \(2\) & \(1+\langle N\rangle\) (Bose--Einstein) & 混沌光 \\
压缩真空 & \(3+\frac{1}{\langle N\rangle}\) & \(2\langle N\rangle+1\) & 光学参量振荡器低于阈值 \\
\hline
\end{tabular}
\end{center}

推导要点。对 Fock 态 \(|n\rangle\)，\(\langle N(N-1)\rangle=n(n-1)\)，故 \(g^{(2)}(0)=1-1/n\)；单光子 \(n=1\) 给出 \(g^{(2)}(0)=0\)，这是单光子源的判据。对相干态，Poisson 分布使 \(\langle N(N-1)\rangle=\langle N\rangle^2\)，故 \(g^{(2)}(0)=1\)。对热光场，Bose--Einstein 分布 \(P_n=\bar n^n/(1+\bar n)^{n+1}\) 给出 \(\langle N(N-1)\rangle=2\bar n^2\)，故 \(g^{(2)}(0)=2\)（光子 bunching）。对压缩真空，只含偶数光子使 \(\langle N(N-1)\rangle=3\sinh^2 r\cosh^2 r\)，故 \(g^{(2)}(0)=3+\coth^2 r>3\)（强 bunching）。

\begin{exbox}[单光子源的 \(g^{(2)}(0)\) 测量]
量子点单光子源在 4\,K、920\,nm 发射，HBT 测得 \(g^{(2)}(0)=0.02\)（接近理想值 0）。残余非零来自多激基态再激发。对比：衰减激光 \(\langle N\rangle=0.01\) 给出 \(g^{(2)}(0)=1\)，无法用于量子密钥分发（BB84 协议要求 \(g^{(2)}(0)<0.5\) 以保证安全性）。
\end{exbox}

理想激光器输出近似相干态。He--Ne 激光器（632.8\,nm，1\,mW）的参数：
\[
\langle N\rangle=\frac{P}{\hbar\omega}\cdot\tau_{\rm det}
\]
对 \(\tau_{\rm det}=1\,\mathrm{ms}\) 探测窗口，\(\langle N\rangle\approx3.2\times10^{12}\)，光子数的相对涨落 \(\Delta N/\langle N\rangle=1/\sqrt{\langle N\rangle}\approx5.6\times10^{-7}\)。

实际激光的 \(g^{(2)}(0)\) 略大于 1（excess noise），取决于工作区间：远高于阈值时 \(g^{(2)}(0)\to1\)（相干态极限），阈值附近 \(g^{(2)}(0)\approx2\)（类热光场）。

\begin{exbox}[相干态的 Poisson 统计验证]
用 HBT 测量衰减脉冲激光（\(\langle N\rangle=5\)），符合计数分布拟合 Poisson \(P_n=e^{-5}5^n/n!\)：\(P_0=0.0067\)，\(P_5=0.1755\)，\(P_{10}=0.0181\)。与热光场的 Bose--Einstein 分布 \(P_n=5^n/6^{n+1}\)（\(P_0=0.167\)，长尾）形成鲜明对比。
\end{exbox}

相干态的实验产生依赖激光器工作于远高于阈值区间。压缩态则需非线性光学过程，主要方案是参量下转换（自发参量下转换），探测则依赖 balanced 零差技术。

参量下转换 (自发参量下转换)。二阶非线性晶体（\(\chi^{(2)}\)）在简并参量下转换配置下实现压缩算符 \(S(\zeta)\)。用 Nd:YAG 泵浦光（532\,nm）注入 PPKTP 晶体，信号场频率为 1064\,nm。光学参量振荡器低于阈值时输出压缩真空，压缩参量
\[
r=\frac{\chi^{(2)}|\mathcal E_p|L}{2c}\sqrt{\frac{\hbar\omega_s}{2\epsilon_0 V_{\rm mode}}}
\]
其中 \(\mathcal E_p\) 是泵浦场振幅，\(L\) 是晶体长度，\(V_{\rm mode}\) 是模体积。

\begin{center}
\begin{tabular}{lcccc}
\hline
实验 & 年份 & 压缩 dB & \(r\) & \(\langle N\rangle=\sinh^2 r\) \\
\hline
Wu \& Kimble (1986) & 1986 & 3.0 & 0.35 & 0.13 \\
LIGO 注入压缩 & 2013 & 6.0 & 0.69 & 0.57 \\
UHF OPO (2016) & 2016 & 15.0 & 1.73 & 8.7 \\
Takinoue (2020) & 2020 & 20.0 & 2.30 & 19.5 \\
\hline
\end{tabular}
\end{center}

压缩 dB 定义为 \(-10\log_{10}(e^{-2r})\)，即 \(20r\log_{10}e\approx8.69\,r\)。15\,dB 压缩意味着噪声功率比真空涨落低 \(10^{1.5}\approx31.6\) 倍。

Balanced 零差探测。压缩态的噪声压缩不在光子数分布中直接可见（压缩真空的光子数只含偶数项），而需通过 balanced 零差探测测量正交分量 \(X_\theta=(ae^{-i\theta}+a^\dagger e^{i\theta})/\sqrt{2}\) 的方差。本地振荡光 (LO) 与信号在 50/50 分束器上混合，两路光电流差
\[
\hat i_-\propto|\alpha_{\rm LO}|\,\hat X_\theta
\]
直接给出正交分量的概率分布。对压缩真空 \(S(\zeta)|0\rangle\)（\(\zeta=r\)，\(\phi=0\)），\(X_0\) 的方差为 \(\frac12 e^{-2r}\)，\(X_{\pi/2}\) 的方差为 \(\frac12 e^{2r}\)。

\begin{exbox}[LIGO 的压缩光注入]
Advanced LIGO 在 2013 年注入 6\,dB 压缩真空，将量子噪声降低约 2 倍。信号频带 10--100\,Hz 中，辐射压力噪声（低频）和散粒噪声（高频）的交叉频率 \(f_c\approx100\,\mathrm{Hz}\)。注入压缩后，高频散粒噪声降低 \(\sim e^{-r}\approx0.50\)，低频辐射压力噪声相应增加 \(e^r\approx2.0\)——但通过滤波器 cavity 将压缩角旋转至频带依赖，可同时降低两个频带（频率依赖压缩, FDC），2020 年 O3 观测运行中实现了 \(\sim3\) 倍的灵敏度提升。
\end{exbox}

Green 并矢\eqnrefs{eq:spontaneous-green-dyadic}不仅适用于自由空间，在腔中其虚部在共振频率处急剧增强，给出 Purcell 效应——自发辐射率被腔共振放大。

\eqnrefs{eq:spontaneous-green-dyadic} 的 Green 并矢在单模腔中给出 Purcell 增强。腔的品质因子 \(Q=\omega_c/\kappa\) 和模体积 \(V_{\rm mode}\) 决定 Purcell 因子
\[
F_P=\frac{\Gamma_{\rm cavity}}{\Gamma_0}
=\frac{3}{4\pi^2}\frac{Q}{V_{\rm mode}/\lambda^3}
\]
其中 \(\lambda=c/\omega_c\) 是共振波长。

\begin{center}
\begin{tabular}{lcccc}
\hline
体系 & \(Q\) & \(V_{\rm mode}/\lambda^3\) & \(F_P\) & \(\Gamma_0/2\pi\) (MHz) \\
\hline
Fabry--P\'erot (Cs) & \(10^5\) & \(10^3\) & \(\sim 8\) & 5.2 \\
微环腔 (InGaAs QD) & \(10^4\) & 0.5 & \(\sim 1500\) & 1.0 \\
光子晶体纳米腔 & \(10^5\) & 0.3 & \(\sim 8000\) & 0.5 \\
超导电路 QED & \(10^6\) & \(10^{-3}\) & \(\sim 8\times10^7\) & 5.0 \\
\hline
\end{tabular}
\end{center}

\begin{exbox}[电路 QED 的强耦合]
超导传输线腔耦合到 transmon 量子比特：腔频率 \(\omega_c/2\pi=7.0\,\mathrm{GHz}\)，比特频率 \(\omega_q/2\pi=6.8\,\mathrm{GHz}\)，耦合强度 \(g/2\pi=100\,\mathrm{MHz}\)，腔衰减率 \(\kappa/2\pi=1\,\mathrm{MHz}\)，比特弛豫率 \(\gamma/2\pi=2\,\mathrm{MHz}\)。协同参数 \(C=4g^2/(\kappa\gamma)=10^4\gg1\)，处于强耦合区。Purcell 增强使比特弛豫率在腔共振时增加 \(F_P\approx4g^2/(\kappa^2)\approx4\times10^4\) 倍——这是电路 QED 中读出比特态的物理基础。
\end{exbox}

前述压缩态均为单模压缩。非简并参量下转换产生的是双模压缩态，信号与闲置模式光子数完全关联，构成连续变量量子纠缠的 EPR 态实现。

非简并自发参量下转换产生双模压缩态
\[
|\mathrm{TMSV}\rangle=S_2(\zeta)|0,0\rangle
=\frac{1}{\cosh r}\sum_{n=0}^\infty(-e^{i\phi}\tanh r)^n|n\rangle_s|n\rangle_i
\]
信号 (signal) 和闲置 (idler) 模式光子数完全关联：\(\langle N_s\rangle=\langle N_i\rangle=\sinh^2 r\)，\(\Delta(N_s-N_i)^2=0\)。这正是连续变量量子纠缠的 EPR 态实现。

\begin{exbox}[EPR 纠缠光的 Reid 判据]
双模压缩态的 Reid 乘积
\[
\Delta(X_s-X_i)\,\Delta(P_s+P_i)=e^{-2r}
\]
对 \(r>0\) 严格小于真空极限 1，验证 EPR 纠缠。2000 年 Bowen 实验用光学参量振荡器产生 \(r\approx0.6\)（3\,dB）双模压缩，测得 Reid 乘积 \(0.64<1\)，首次实验验证连续变量 EPR 纠缠。
\end{exbox}

\begin{proof}[解]

Wigner 函数与特征函数。单模 Wigner 函数定义为
\begin{equation}
W(\alpha)=\frac{1}{\pi}\Tr\!\bigl[D^\dagger(\alpha)\,\rho\,D(\alpha)\,(-1)^{\hat n}\bigr],
\qquad
\alpha=\frac{x_1+\ii x_2}{\sqrt2},
\label{eq:wigner-definition}
\end{equation}
其中 \(D(\alpha)=\exp(\alpha a^\dagger-\alpha^*a)\) 为位移算符，\((-1)^{\hat n}\) 为宇称算符。等价地通过特征函数
\begin{equation}
\chi(\xi)=\Tr\!\bigl[\rho\,D(\xi)\bigr],
\qquad
W(\alpha)=\frac{1}{\pi^2}\int\!\dd^2\xi\;\ee^{\alpha\xi^*-\alpha^*\xi}\,\chi(\xi).
\label{eq:wigner-characteristic}
\end{equation}
真空态 \(\rho_0=|0\rangle\langle0|\) 的特征函数 \(\chi_0(\xi)=\langle0|D(\xi)|0\rangle=\ee^{-|\xi|^2/2}\)，代入\eqnrefs{eq:wigner-characteristic} 完成高斯积分得
\begin{equation}
W_0(\alpha)=\frac{2}{\pi}\,\ee^{-2|\alpha|^2}
=\frac{2}{\pi}\,\ee^{-(x_1^2+x_2^2)}.
\label{eq:vacuum-wigner}
\end{equation}

压缩态特征函数。其次考虑单模压缩态 \(\rho_{\rm sq}=S(\zeta)|0\rangle\langle0|S^\dagger(\zeta)\)，\(\zeta=r\ee^{\ii\phi}\)。特征函数
\begin{equation*}
\chi_{\rm sq}(\xi)=\Tr\!\bigl[S(\zeta)|0\rangle\langle0|S^\dagger(\zeta)\,D(\xi)\bigr]
=\langle0|S^\dagger(\zeta)D(\xi)S(\zeta)|0\rangle.
\end{equation*}
利用 Bogoliubov 变换（即\eqnrefs{eq:displacement-normal-order} 的算符形式）
\begin{equation}
S^\dagger(\zeta)\,D(\xi)\,S(\zeta)
=D\!\bigl(\xi\cosh r+\xi^*\ee^{\ii\phi}\sinh r\bigr),
\label{eq:displacement-bogoliubov}
\end{equation}
此式由本章约定 \(S^\dagger a S=a\cosh r-a^\dagger\ee^{\ii\phi}\sinh r\) 代入 \(D(\xi)=\exp(\xi a^\dagger-\xi^*a)\) 得到；位移参数中的正号来自同时变换 \(a\) 与 \(a^\dagger\)。故
\begin{equation*}
\chi_{\rm sq}(\xi)=\chi_0\!\bigl(\xi\cosh r+\xi^*\ee^{\ii\phi}\sinh r\bigr)
=\exp\!\Bigl[-\frac12\bigl|\xi\cosh r+\xi^*\ee^{\ii\phi}\sinh r\bigr|^2\Bigr].
\end{equation*}

\end{proof}

\begin{proof}[解]
取 \(\phi=0\)，此时 \(\xi=\xi_1+\ii\xi_2\)，
\begin{equation*}
\bigl|\xi\cosh r+\xi^*\sinh r\bigr|^2
=(\xi_1\ee^r)^2+(\xi_2\ee^{-r})^2
=\ee^{2r}\xi_1^2+\ee^{-2r}\xi_2^2.
\end{equation*}
代入\eqnrefs{eq:wigner-characteristic}，Fourier 逆变换为两个独立高斯积分：
\begin{equation}
W_{\rm sq}(\alpha)
=\frac{2}{\pi}\,\exp\!\Bigl[-x_1^2\ee^{2r}-x_2^2\ee^{-2r}\Bigr],
\label{eq:squeezed-wigner-function}
\end{equation}
即\eqnrefs{eq:squeezed-wigner-function}。相空间中 Wigner 函数是椭圆高斯：沿 \(x_1\) 方向压缩、沿 \(x_2\) 方向展宽，而协方差行列式保持不变。

注意 \(\dd^2\alpha=\dd x_1\dd x_2/2\)。因此归一化边缘分布为
\begin{equation*}
P(x_1)=\frac12\int_{-\infty}^{+\infty}\!\dd x_2\;W_{\rm sq}(\alpha)
=\frac{\ee^r}{\sqrt\pi}\,\ee^{-x_1^2\ee^{2r}},
\end{equation*}
\begin{equation*}
P(x_2)=\frac12\int_{-\infty}^{+\infty}\!\dd x_1\;W_{\rm sq}(\alpha)
=\frac{\ee^{-r}}{\sqrt\pi}\,\ee^{-x_2^2\ee^{-2r}}.
\end{equation*}
方差
\begin{equation}
\Delta X_1=\frac{\ee^{-r}}{\sqrt2},
\qquad
\Delta X_2=\frac{\ee^r}{\sqrt2},
\label{eq:squeezed-quadrature-variance}
\end{equation}
即\eqnrefs{eq:squeezed-quadrature-variance}。\(r>0\) 时 \(X_1\) 压缩、\(X_2\) 反压缩，与 Wigner 椭圆的轴长一致。

不确定性原理饱和。最后计算乘积
\begin{equation*}
\Delta X_1\,\Delta X_2
=\frac{\ee^{-r}}{\sqrt2}\cdot\frac{\ee^r}{\sqrt2}
=\frac12,
\end{equation*}
恰为 Heisenberg 下界 \(\Delta X_1\,\Delta X_2\ge1/2\) 的等号。真空态 \(r=0\) 给出对称噪声；\(r>0\) 只是在两个共轭分量之间重新分配噪声。

\end{proof}

压缩态的量子信息性质可由正交分量方差和高斯熵继续分析。

标准量子极限与压缩超越。相干态 \(|\alpha\rangle\) 的正交分量 \(\Delta X=\Delta P=1/\sqrt{2}\)（最小不确定乘积 \(\Delta X\Delta P=1/2\)），给出标准量子极限（SQL）。压缩态 \(S(r)|\alpha\rangle\)（\(S(r)=\exp[\frac{r}{2}(a^2-a^{\dagger2})]\)）满足 \(\Delta X=\ee^{-r}/\sqrt{2}\)，\(\Delta P=\ee^r/\sqrt{2}\)——一个分量压缩，另一个反压缩，乘积仍 \(1/2\)。压缩 \(r\) 使 \(\Delta X\) 低于 SQL，在引力波探测（LIGO 注入压缩光提升灵敏度）和精密测量中直接应用。

Holevo 界与高斯信道容量。连续变量量子信道的经典信息容量由 Holevo 界给出：\(C\le\chi=S(\rho_{\rm out})-\int p(x)S(\rho_x)\dd x\)（\(S\) von Neumann 熵）。对高斯信道（保高斯态的 CP 映射），Holevo 界可精确计算：\(C=g(\nu_1)+g(\nu_2)-g(\nu_3)-g(\nu_4)\)（\(g(x)=(x+1)\log(x+1)-x\log x\)，\(\{\nu_i\}\) 是信道参数的辛本征值）。压缩态优化 Holevo 界——在能量约束下，压缩态（而非相干态）达到高斯信道容量，是连续变量量子通信的最优编码。

纠缠辅助容量与超加性。纠缠辅助容量 \(C_E\)（发送方与接收方共享纠缠时的容量）可超过未辅助容量 \(C\)（超加性 \(C_E\ge C\)）。对高斯信道，\(C_E=g(N_S)-g(N_N)\)（\(N_S\) 信号光子数，\(N_N\) 噪声光子数），而 \(C\) 需要优化输入态。对纯损耗信道，\(C=C_E=\log(1+\eta N_S/(1-\eta))\)（\(\eta\) 透射率），纠缠不增加容量；对噪声信道，\(C_E>C\)，纠缠提供真实增益——这是量子通信中纠缠资源价值的定量体现。

多模压缩与连续变量纠缠。双模压缩态 \(S_2(r)|0,0\rangle=\frac{1}{\cosh r}\sum_n(\tanh r)^n|n,n\rangle\) 是连续变量纠缠的典型态。其 Schmidt 分解 \(\sum_n\lambda_n|n\rangle_A|n\rangle_B\)（\(\lambda_n=(\tanh r)^n/\cosh r\)）给出纠缠熵 \(E=-\sum|\lambda_n|^2\log|\lambda_n|^2=g(\sinh^2 r)\)。双模压缩态在连续变量量子密钥分发（CV-QKD）中用作纠缠源，安全性的 Shor--Preskill 证明在连续变量中通过 Holevo 界给出窃听者信息上界 \(I_{\rm Eve}\le\chi\)——压缩态参数 \(r\) 直接决定密钥率 \(K=\beta I_{AB}-\chi\)（\(\beta\) 重复率，\(I_{AB}\) Alice--Bob 互信息），是连续变量量子通信的物理基础。
